\chapter{System Design and Methodology}

\section{Introduction}
%TODO%

\section{High-Level System Architecture}
\label{sec:architecture}
% ------------------------------------------------------------
 
The proposed system is designed as a distributed robotic architecture
structured across three main layers: the \emph{embedded device layer}
(Raspberry Pi and Arduino), the \emph{server layer}, and the
\emph{user interface layer} (web dashboard). Together these layers
form a real-time, event-driven pipeline that integrates perception,
communication, and physical actuation into a cohesive autonomous
system.
 
\bigskip
\noindent\textbf{[DIAGRAM 1 — System Architecture Overview]}\\
\textit{A three-tier block diagram showing the Embedded Layer
(Raspberry Pi + Arduino), the Server Layer (backend + inference
engine), and the UI Layer (dashboard), with labelled arrows indicating
WebSocket and serial communication channels between them. Place after
this paragraph.}
\bigskip
 
% ............................................................
\subsection{System Startup and Initialisation}
% ............................................................
 
On power-up, the Raspberry Pi executes an initialisation sequence
before entering normal operation. This sequence serves two purposes:
it verifies that the network environment is reachable, and it
authenticates any remote client wishing to interact with the
system~\cite{heath2002embedded}.
 
The initialisation procedure follows these steps:
 
\begin{enumerate}
    \item \textbf{Network check.} The system tests for an active
          network connection. If none is available, it cannot
          establish the server link required for remote monitoring
          and fallback inference.
 
    \item \textbf{Credential verification.} Authentication credentials
          for remote access are validated. If credentials do not
          match, the system enters a \emph{protected provisioning
          mode}: it suspends external control and waits for a
          successful authenticated handshake before proceeding.
\end{enumerate}
 
This two-step gate ensures that the system remains inoperable to
unauthorised parties even if the hardware is physically accessible.
 
% ............................................................
\subsection{Communication Initialisation}
% ............................................................
 
Once the initialisation guard passes, two independent communication
channels are brought up concurrently:
 
\begin{itemize}
    \item \textbf{WebSocket connection} between the Raspberry Pi and
          the backend server. WebSockets provide a persistent,
          full-duplex TCP channel that allows either side to push
          messages without polling~\cite{fette2011websocket}. This
          channel carries high-level control commands, video frame
          streams, and telemetry data.
 
    \item \textbf{Serial (UART) connection} between the Raspberry Pi
          and the Arduino microcontroller. Serial communication
          provides deterministic, low-latency delivery of binary
          command packets to hardware peripherals — suitable for
          real-time actuator control where the overhead of a
          network stack would be unnecessary~\cite{axelson2007serial}.
\end{itemize}
 
The two channels are kept deliberately separate: the WebSocket handles
supervisory and data-plane traffic while the serial link handles
time-critical hardware commands. This separation prevents high-level
communication delays from affecting actuator response times.
 
\bigskip
\noindent\textbf{[DIAGRAM 2 — Communication Channels]}\\
\textit{A sequence or swim-lane diagram showing the Raspberry Pi in
the centre, with WebSocket traffic flowing left to the server and
serial traffic flowing right to the Arduino. Include example message
types on each arrow (e.g., ``frame stream'', ``pump ON''). Place
after this subsection.}
\bigskip
 
% ............................................................
\subsection{Multi-Threaded Execution Architecture}
% ............................................................
 
After initialisation, the Raspberry Pi enters its main execution
phase. To prevent any single task from blocking the others, the
system adopts a \emph{multi-threaded} design in which four concurrent
threads run in parallel:
 
\begin{enumerate}
    \item \textbf{Camera capture thread.} Continuously acquires frames
          from the camera module at the configured frame rate and
          places them in a shared buffer for the inference thread to
          consume.
 
    \item \textbf{AI inference thread.} Dequeues frames from the
          buffer and runs the detection pipeline (see
          Section~\ref{sec:inference_pipeline}). Publishes detection
          events to a shared event queue.
 
    \item \textbf{Communication thread.} Manages the WebSocket
          connection (sending frames, receiving commands) and the
          Arduino serial link (forwarding actuation commands).
 
    \item \textbf{State machine thread.} Reads the event queue and
          advances the Finite State Machine (FSM) described in
          Section~\ref{sec:fsm}, translating system state transitions
          into concrete actions.
\end{enumerate}
 
Thread-safe queues and mutexes mediate access to shared resources,
ensuring that frame buffers and event queues are read and written
without race conditions~\cite{silberschatz2018os}.
 
\bigskip
\noindent\textbf{[DIAGRAM 3 — Thread Architecture]}\\
\textit{A concurrent thread diagram (four horizontal swim lanes)
showing the Camera, Inference, Communication, and FSM threads,
with shared data structures (frame buffer, event queue) indicated
as boxes between the relevant lanes. Place after this subsection.}
\bigskip
 
% ............................................................
\subsection{Perception and Inference Pipeline}
\label{sec:inference_pipeline}
% ............................................................
 
The core perception task — detecting pests or target objects in the
camera feed — is implemented as a two-mode inference pipeline with an
automatic fallback mechanism:
 
\begin{enumerate}
    \item \textbf{Local inference (primary mode).} If a compatible
          optimised model (e.g., a quantised TFLite or ONNX file) is
          present on the Raspberry Pi, inference is performed
          on-device. This is the preferred path: it minimises latency,
          avoids network dependency, and keeps the detection loop
          self-contained.
 
    \item \textbf{Remote inference (fallback mode).} If no local model
          is available or the device lacks sufficient resources,
          captured frames are streamed to the backend server over
          the WebSocket channel. The server performs inference and
          returns detection results to the Pi.
\end{enumerate}
 
This dual-mode design improves system robustness: the robot continues
to function even under degraded computational conditions, and the
same codebase supports both embedded and server-side execution without
modification.
 
\bigskip
\noindent\textbf{[DIAGRAM 4 — Inference Pipeline Flowchart]}\\
\textit{A flowchart beginning at ``New Frame Captured'', branching on
``Local model available?'' — Yes path leads to on-device inference;
No path leads to WebSocket transmission and server inference. Both
paths converge at ``Detection Result Available''. Place after this
subsection.}
\bigskip
 
% ............................................................
\subsection{Event-Driven Actuation}
% ............................................................
 
When the inference pipeline produces a positive detection, the system
responds through an event-driven actuation chain that cleanly
separates the perception and control concerns:
 
\begin{enumerate}
    \item \textbf{Detection event raised.} The inference thread pushes
          a detection event (object class, bounding box, confidence
          score) onto the shared event queue.
 
    \item \textbf{Control signal generated.} The FSM consumes the
          event, transitions to the \emph{Actuation} state, and
          generates the appropriate command (e.g., activate pump).
 
    \item \textbf{Command forwarded to Arduino.} The communication
          thread serialises the command and transmits it over the
          UART link to the Arduino.
 
    \item \textbf{Physical actuation.} The Arduino executes the
          command by driving the relevant peripheral — for example,
          triggering a relay to activate the pesticide pump.
\end{enumerate}
 
Decoupling perception from actuation in this way means that the
detection logic can be modified or replaced independently of the
hardware control logic, improving modularity and simplifying
testing~\cite{brooks1991intelligence}.
 
\bigskip
\noindent\textbf{[DIAGRAM 5 — Actuation Event Chain]}\\
\textit{A linear pipeline diagram: Inference Thread
$\rightarrow$ Event Queue $\rightarrow$ FSM $\rightarrow$
Communication Thread $\rightarrow$ Arduino $\rightarrow$ Pump/Actuator.
Label each arrow with the data passed (event struct, UART command
bytes, relay signal). Place after this subsection.}
\bigskip
 
% ............................................................
\subsection{Server-Side Architecture}
% ............................................................
 
The backend server functions as both a \emph{computational extension}
of the Raspberry Pi and a \emph{supervisory node} for the entire
system. Its responsibilities are:
 
\begin{itemize}
    \item \textbf{Frame stream reception.} Accepts compressed video
          frames from the Pi over WebSockets and, when acting in
          fallback inference mode, passes them through the detection
          pipeline.
 
    \item \textbf{Lazy model loading.} Detection models are loaded
          into memory only when the first inference request arrives,
          reducing idle memory consumption on the
          server~\cite{gamma1994design}.
 
    \item \textbf{Health monitoring.} The server tracks heartbeat
          signals from the Pi and maintains a live view of system
          state, flagging connectivity loss or processing anomalies.
 
    \item \textbf{Remote control relay.} Commands issued by the user
          through the dashboard are validated and forwarded to the Pi
          over the persistent WebSocket connection.
\end{itemize}
 
% ............................................................
\subsection{User Dashboard}
% ............................................................
 
The web-based dashboard provides the human operator with a graphical
interface for monitoring and controlling the system without requiring
direct access to the embedded hardware. Through the dashboard, an
operator can:
 
\begin{itemize}
    \item Start or stop the live camera stream.
    \item Issue manual control commands to the robot.
    \item Trigger a remote reboot or system reset of the Raspberry Pi.
    \item Monitor real-time system health metrics (CPU load,
          temperature, connection status, detection counts).
\end{itemize}
 
The dashboard communicates exclusively with the backend server;
it has no direct connection to the embedded hardware, which keeps
the interface layer isolated from hardware-specific
details~\cite{gamma1994design}.
 
\bigskip
\noindent\textbf{[DIAGRAM 6 — Dashboard UI Mockup / Screenshot]}\\
\textit{A screenshot or wireframe of the dashboard showing the live
stream panel, status indicators, and control buttons. Place here if
the UI is implemented; otherwise a wireframe is sufficient.}
\bigskip
 
% ............................................................
\subsection{Finite State Machine (FSM) for System Control}
\label{sec:fsm}
% ............................................................
 
The overall behaviour of the Raspberry Pi is governed by a
\emph{Finite State Machine} (FSM)~\cite{hopcroft2006automata}. An FSM
is a computational model that defines a fixed set of states and the
conditions under which the system transitions between them. Using an
FSM makes system behaviour explicit, predictable, and easy to
validate — properties that are essential in a safety-relevant
agricultural robotics context.
 
The FSM defines five states:
 
\begin{itemize}
    \item \textbf{Idle.} The system is initialised and connected but
          not actively processing. Waiting for a user command or
          scheduled trigger.
 
    \item \textbf{Streaming.} The camera capture and communication
          threads are active; frames are being acquired and forwarded.
          No detection is running.
 
    \item \textbf{Detection.} The inference thread is active;
          frames are being analysed for target objects.
 
    \item \textbf{Actuation.} A detection event has been confirmed;
          the system is executing the corresponding physical action
          (e.g., spraying). The camera and inference threads continue
          running.
 
    \item \textbf{Error / Recovery.} An anomaly has been detected
          (network loss, hardware fault, inference failure). The system
          suspends normal operation, logs the fault, and attempts
          recovery before returning to Idle.
\end{itemize}
 
\bigskip
\noindent\textbf{[DIAGRAM 7 — FSM State Diagram]}\\
\textit{A standard UML state diagram with the five states as rounded
rectangles and labelled directed arrows for each transition (e.g.,
``user starts stream'' Idle $\rightarrow$ Streaming; ``detection
confirmed'' Detection $\rightarrow$ Actuation; ``fault raised'' any
state $\rightarrow$ Error). This is one of the most important diagrams
in the chapter — place immediately after this subsection.}
\bigskip

\section{Hardware Architecture}
\label{sec:hardware}
% ------------------------------------------------------------
 
This section describes the hardware design of the system, covering the
selection and role of each component, the power distribution strategy,
and the inter-component communication scheme. The design follows a
\emph{modular, distributed} approach: a high-level computing unit
handles perception and decision-making, while a dedicated
microcontroller manages real-time actuation. This separation reflects
a well-established principle in embedded robotics — keeping
computationally intensive tasks away from time-critical control
loops~\cite{heath2002embedded}.
 
It is important to note upfront that the hardware design described here
represents the \emph{full intended architecture}. The implemented
prototype integrates only the Raspberry Pi, Arduino, and camera module;
the motor and pump subsystems remain part of the conceptual design and
are documented here to provide a complete picture of the target
platform and to guide future implementation.
 
\bigskip
\noindent\textbf{[DIAGRAM 1 — Hardware Block Diagram]}\\
\textit{A block diagram showing all hardware components and their
interconnections: Raspberry Pi (centre) connected to the Pi Camera
(top), Arduino (right via USB), battery and buck converter (bottom
left), and motor drivers + motors/pump (bottom right). Shade the
implemented components differently from the conceptual ones to make
the prototype boundary immediately visible. Place after this
introduction paragraph.}
\bigskip
 
% ............................................................
\subsection{Processing and Vision Unit — Raspberry Pi}
% ............................................................
 
The Raspberry Pi serves as the central processing unit of the system.
It is responsible for three core functions: running the computer vision
and object detection pipeline, executing the high-level system logic
and Finite State Machine (see Section~\ref{sec:fsm}), and coordinating
all communication between layers.
 
The Raspberry Pi was selected for this role because it combines a
multi-core ARM CPU with a full Linux environment, native Python
support, and hardware interfaces (USB, GPIO, CSI camera port) that
are sufficient for prototyping AI-enabled robotic systems while
remaining within the power and cost constraints of an embedded
platform~\cite{upton2016raspberry}.
 
\paragraph{Camera module.}
A Raspberry Pi Camera Module is connected directly to the board via
the dedicated CSI (Camera Serial Interface) ribbon connector. The CSI
interface offloads image transfer from the USB bus, providing a
dedicated high-bandwidth, low-latency path for continuous frame
acquisition~\cite{picamera2023}. The camera feeds raw frames into the
inference pipeline described in Section~\ref{sec:inference_pipeline}.
 
% ............................................................
\subsection{Microcontroller Control Unit — Arduino}
% ............................................................
 
An Arduino microcontroller acts as the low-level control unit,
responsible for executing hardware commands with deterministic timing.
Microcontrollers are preferred over the Raspberry Pi for this role
because they run bare-metal or RTOS firmware with no operating system
scheduler overhead, guaranteeing response times in the order of
microseconds rather than the milliseconds typical of a
Linux-scheduled process~\cite{heath2002embedded}.
 
The Arduino is connected to the Raspberry Pi over a \textbf{USB serial
(UART)} link. This connection serves two purposes simultaneously: it
supplies power to the Arduino from the Raspberry Pi's USB port, and it
provides the communication channel through which the Pi sends command
packets and the Arduino returns status acknowledgements.
 
Upon receiving a command, the Arduino decodes it and drives the
appropriate GPIO pins to control motors, relays, or other peripherals
directly, without further involvement from the Raspberry Pi. This
ensures that time-critical signals — such as a precise PWM waveform
for motor speed control — are generated by hardware dedicated to
that task alone.
 
\bigskip
\noindent\textbf{[DIAGRAM 2 — Raspberry Pi $\leftrightarrow$ Arduino
Communication Detail]}\\
\textit{A zoomed-in diagram showing the USB serial link between the
two boards, with example command/acknowledgement byte sequences
annotated on the arrows (e.g., \texttt{CMD:PUMP\_ON} $\rightarrow$,
\texttt{ACK:OK} $\leftarrow$). Place after this subsection.}
\bigskip
 
% ............................................................
\subsection{Actuation System — Motors and Pump Mechanism}
% ............................................................
 
The actuation layer is designed around DC motors interfaced through
dedicated motor driver circuits. DC motors are driven by voltages and
currents far beyond what a microcontroller GPIO pin can supply
directly; motor drivers act as power amplifiers, accepting a
low-power control signal from the Arduino and switching the full motor
supply voltage accordingly~\cite{monk2016electronics}.
 
The intended actuation configuration comprises two subsystems:
 
\begin{itemize}
    \item \textbf{Mobility subsystem.} Four DC motors arranged in a
          \emph{differential drive} configuration — two on each side
          of the chassis. Differential drive steers the robot by
          varying the relative speed of the left and right motor
          pairs; turning is achieved without a dedicated steering
          mechanism, simplifying both hardware and
          control~\cite{siegwart2011mobile}.
 
    \item \textbf{Spraying subsystem.} Two DC motors dedicated to
          driving the pump mechanism used for pesticide delivery.
          These motors are controlled independently of the mobility
          subsystem, allowing spraying to occur while the robot is
          stationary or in motion.
\end{itemize}
 
Each motor group is controlled through its own motor driver circuit,
which receives PWM (Pulse-Width Modulation) signals from the Arduino.
By varying the PWM duty cycle, the Arduino adjusts motor speed
continuously without the energy loss associated with resistive
control~\cite{monk2016electronics}.
 
\bigskip
\noindent\textbf{[DIAGRAM 3 — Actuation Subsystem Layout]}\\
\textit{A top-down schematic of the robot chassis showing the
four drive motors (front-left, front-right, rear-left, rear-right),
the two pump motors, and the motor driver modules that connect them
to the Arduino. A separate inset can illustrate the differential
drive steering principle (left motors faster $\Rightarrow$ turn
right). Place after this subsection.}
\bigskip
 
% ............................................................
\subsection{Power Distribution System}
% ............................................................
 
Reliable power delivery is a critical design concern in mobile robotic
systems. High-current actuators can cause sudden voltage drops on a
shared supply rail, which may reset or corrupt the operation of
sensitive digital electronics~\cite{horowitz2015art}. To prevent this,
the power system is partitioned into two isolated domains:
 
\paragraph{Computational domain.}
The Raspberry Pi is powered from the main battery through a
\textbf{5\,V DC-DC buck converter}. A buck converter steps down the
battery voltage efficiently (typically 85–95\% efficiency) and
regulates the output to a stable 5\,V regardless of load
variations~\cite{rashid2017power}. The Arduino derives its power
directly from the Raspberry Pi's USB port (also 5\,V), so both
computing units share the regulated supply without requiring a second
converter.
 
\paragraph{Actuation domain.}
The DC motors and pump mechanism are powered directly from the main
battery, bypassing the buck converter. This provides the motors with
the full battery voltage and ensures that current spikes during motor
start-up or stall do not propagate into the regulated 5\,V rail
feeding the computing units.
 
Table~\ref{tab:power} summarises the power domains of each component.
 
\begin{table}[h]
\centering
\caption{Power supply domain for each system component.}
\label{tab:power}
\begin{tabular}{llcc}
\hline
\textbf{Component} & \textbf{Supply Path} & \textbf{Voltage} & \textbf{Domain} \\
\hline
Raspberry Pi     & Battery $\to$ Buck converter & 5\,V  & Computational \\
Arduino          & Raspberry Pi USB port         & 5\,V  & Computational \\
Pi Camera Module & Raspberry Pi CSI port         & 3.3\,V & Computational \\
DC Drive Motors  & Battery (direct)              & Battery & Actuation \\
Pump Motors      & Battery (direct)              & Battery & Actuation \\
\hline
\end{tabular}
\end{table}
 
\bigskip
\noindent\textbf{[DIAGRAM 4 — Power Distribution Schematic]}\\
\textit{A power-flow diagram starting from the battery, branching
into the buck converter path (to Pi and Arduino) and the direct path
(to motor drivers and motors). Use colour coding to distinguish the
two power domains clearly. This diagram is important for
understanding the isolation strategy — place it immediately after
Table~\ref{tab:power}.}
\bigskip
 
% ............................................................
\subsection{Inter-Component Communication Summary}
% ............................................................
 
Table~\ref{tab:comms} summarises the communication interfaces used
between hardware components. Two physical interfaces are used in the
system, each chosen to match the bandwidth and latency requirements
of the data it carries.
 
\begin{table}[h]
\centering
\caption{Hardware communication interfaces.}
\label{tab:comms}
\begin{tabular}{llll}
\hline
\textbf{Link} & \textbf{Interface} & \textbf{Direction} & \textbf{Data} \\
\hline
Pi $\leftrightarrow$ Server  & WebSocket (TCP/IP, Wi-Fi) & Bidirectional & Frames, commands, telemetry \\
Pi $\leftrightarrow$ Arduino & USB Serial (UART)         & Bidirectional & Control commands, ACKs \\
Pi $\leftrightarrow$ Camera  & CSI ribbon cable          & Unidirectional & Raw image frames \\
Arduino $\leftrightarrow$ Motors & PWM via motor driver & Unidirectional & Speed / direction signals \\
\hline
\end{tabular}
\end{table}
 
The \textbf{CSI} interface between the Pi and the camera is
unidirectional and hardware-accelerated, offloading frame transfer
from the main CPU bus. The \textbf{USB serial} link between the Pi
and Arduino is the sole path for all actuation commands; keeping it
as a single, auditable channel simplifies debugging and ensures that
every hardware action can be traced to a logged command.
 
% ............................................................
\subsection{Prototype Scope vs.\ Full Design}
% ............................................................
 
As stated at the outset of this section, the current prototype
implements a subset of the full hardware design. Table~\ref{tab:scope}
maps each hardware subsystem to its implementation status.
 
\begin{table}[h]
\centering
\caption{Implementation status of hardware subsystems.}
\label{tab:scope}
\begin{tabular}{lcc}
\hline
\textbf{Subsystem} & \textbf{In Full Design} & \textbf{In Prototype} \\
\hline
Raspberry Pi (SBC)         & \checkmark & \checkmark \\
Pi Camera Module           & \checkmark & \checkmark \\
Arduino microcontroller    & \checkmark & \checkmark \\
DC drive motors + drivers  & \checkmark & $\times$   \\
Pump motors + drivers      & \checkmark & $\times$   \\
Buck converter (5\,V)      & \checkmark & \checkmark \\
Battery power supply       & \checkmark & $\times$   \\
\hline
\end{tabular}
\end{table}
 
This scoping decision reflects the priority of the current project
phase: validating the software architecture, communication pipeline,
and model deployment strategy on real hardware before committing to
the mechanical and electrical integration of the mobility and
spraying subsystems. The modular design ensures that adding the
remaining hardware components in a future iteration requires no
changes to the software or communication architecture described in
Section~\ref{sec:architecture}.


\section{Software Architecture}
\label{sec:software}
% ------------------------------------------------------------
 
The software architecture is designed around a clear separation of
concerns between two independent subsystems: the \emph{Raspberry Pi
software stack}, which handles high-level coordination, perception,
and communication, and the \emph{Arduino firmware}, which executes
time-critical, low-level actuator control. This division mirrors the
hardware partitioning described in Section~\ref{sec:hardware} and
reflects a fundamental principle of embedded systems design — keeping
computationally intensive tasks structurally isolated from real-time
control loops so that neither can interfere with the other's timing
guarantees~\cite{heath2002embedded}.
 
The remainder of this section describes each subsystem in turn,
covering internal module structure, inter-module data flow, and the
safety mechanisms embedded at each layer.
 
\bigskip
\noindent\textbf{[DIAGRAM 1 — Top-Level Software Partitioning]}\\
\textit{A two-column block diagram: left column shows the Raspberry Pi
software stack (five layers stacked vertically); right column shows the
Arduino firmware (four modules). A bidirectional arrow labelled
``UART serial'' connects the two columns at the communication boundary.
Place after this introduction paragraph.}
\bigskip
 
% ------------------------------------------------------------
\subsection{Arduino Firmware Architecture}
% ------------------------------------------------------------
 
The Arduino firmware operates as a \emph{serial-controlled actuator
node}: it exposes a simple command interface over UART and translates
received commands into hardware signals, completely hiding the
electrical details of motor drivers and relay circuits from the host
system. This abstraction means the Raspberry Pi never needs to know
how many PWM pins a motor requires or what logic level activates a
relay — it simply sends a named command and waits for an
acknowledgement.
 
The firmware is organised into four modules, each with a single
well-defined responsibility:
 
\paragraph{1. Actuator Abstraction Layer (AAL).}
The AAL maintains the \emph{logical state} of every actuator in
software — motor speed (as a normalised value), direction (forward /
reverse / brake), and pump status (on / off). All other modules
interact with actuators through the AAL rather than manipulating
hardware registers directly. This indirection means that if a motor
driver is swapped for a different model, only the Hardware Abstraction
Layer needs to change~\cite{gamma1994design}.
 
\paragraph{2. Hardware Abstraction Layer (HAL).}
The HAL translates the logical actuator states maintained by the AAL
into concrete hardware signals. For motors this means computing PWM
duty cycles and writing them to the appropriate timer registers; for
the pump relay it means asserting or deasserting a GPIO pin. The HAL
is the only module that has direct knowledge of pin assignments and
hardware-specific timing requirements~\cite{heath2002embedded}.
 
\paragraph{3. Command Processor.}
The Command Processor reads newline-delimited text frames from the
UART receive buffer, parses and validates them, and dispatches the
corresponding actions to the AAL. Using human-readable, text-based
commands (rather than a binary protocol) simplifies debugging — a
serial monitor can be attached to inspect all traffic without a
protocol decoder. Validation rejects malformed or out-of-range
commands before they reach the AAL, preventing erroneous actuator
states from being applied~\cite{axelson2007serial}.
 
\paragraph{4. Failsafe Monitor.}
The Failsafe Monitor implements a \emph{watchdog-style heartbeat
mechanism}. Whenever any actuator is in an active state, the monitor
checks that a valid command or keep-alive packet has been received
within a configurable timeout window. If communication is interrupted
— due to host software failure, USB disconnection, or physical cable
fault — the monitor immediately commands all actuators to their idle
state and sets an internal error flag. This fail-safe behaviour
prevents the robot from continuing to move or spray uncontrollably in
the absence of a supervising host~\cite{kopetz2011real}.
 
\bigskip
\noindent\textbf{[DIAGRAM 2 — Arduino Firmware Module Diagram]}\\
\textit{Four stacked boxes (Command Processor on top, AAL below it,
HAL below that, physical hardware at the bottom) with a vertical data
flow arrow labelled ``UART RX'' entering the Command Processor from
the left and ``PWM / GPIO'' exiting the HAL to the right. The
Failsafe Monitor sits as a side box with a monitoring arrow pointing
into the AAL. Place after this subsection.}
\bigskip
 
% ------------------------------------------------------------
\subsection{Raspberry Pi Software Architecture}
% ------------------------------------------------------------
 
The Raspberry Pi software stack is responsible for everything above
the hardware serial link: acquiring sensor data, running AI inference,
managing external communication, and making decisions that drive the
rest of the system. The stack follows a \emph{layered,
event-driven architecture}, where each layer has a defined role and
communicates with adjacent layers through well-specified
interfaces~\cite{bass2021software}.
 
The five layers, from lowest to highest level of abstraction, are
described below.
 
\bigskip
\noindent\textbf{[DIAGRAM 3 — Raspberry Pi Five-Layer Stack]}\\
\textit{Five horizontal bands stacked vertically, labelled (bottom to
top): Hardware Layer, State Management Layer, Services Layer,
Application Layer, Utility Layer. Annotate each band with its key
components (e.g., Services Layer: Camera Service, Inference Service,
Communication Service, Health Monitor, Watchdog). Vertical arrows show
upward data flow and downward control flow. Place before the layer
descriptions below.}
\bigskip
 
\paragraph{1. Hardware Layer.}
The Hardware Layer provides software interfaces to physical devices.
It wraps the camera module (via the CSI interface) and the Arduino
serial link (via USB UART) behind consistent Python objects, so that
higher layers acquire a frame or send a command through the same API
regardless of the underlying driver implementation. This isolates
driver-specific code to a single location and allows hardware to be
simulated during unit testing~\cite{gamma1994design}.
 
\paragraph{2. State Management Layer.}
The State Management Layer hosts the central \emph{Finite State
Machine} (FSM) that governs the robot's operational lifecycle (see
Section~\ref{sec:fsm} for the full state diagram). The FSM is the
sole authority on what the system is allowed to do at any given
moment: it gates service start-up, controls state transitions in
response to events from the Services Layer, and triggers the
appropriate shutdown or recovery procedures on fault detection. All
other layers interact with the FSM by posting events; they never
transition states directly~\cite{hopcroft2006automata}.
 
\paragraph{3. Services Layer.}
The Services Layer forms the operational core of the runtime
environment. It comprises five concurrent services, each running in
its own thread:
 
\begin{itemize}
    \item \textbf{Camera Service.} Continuously acquires frames from
          the Hardware Layer and places them on a shared frame queue
          for consumption by the Inference Service.
 
    \item \textbf{Inference Service.} Dequeues frames and runs the
          object detection model. On a positive detection it posts
          a detection event to the FSM event queue. The service
          transparently selects between local and remote inference
          modes as described in Section~\ref{sec:inference_pipeline}.
 
    \item \textbf{Communication Service.} Manages the WebSocket
          connection to the backend server (outgoing frame stream,
          incoming control commands) and the serial connection to the
          Arduino (outgoing actuation commands, incoming status
          acknowledgements).
 
    \item \textbf{Health Monitor.} Periodically samples CPU load,
          memory usage, temperature, and connection status, and
          publishes this telemetry to the server and the local log.
 
    \item \textbf{Watchdog.} Independently verifies that each service
          thread is still alive and processing. If a thread becomes
          unresponsive — due to deadlock, an unhandled exception, or
          resource exhaustion — the Watchdog raises a fault event
          to the FSM, which then initiates a controlled recovery or
          shutdown sequence~\cite{kopetz2011real}.
\end{itemize}
 
\paragraph{4. Application Layer.}
The Application Layer is the entry point of the system. On startup it
reads the configuration, instantiates all services and hardware
drivers, runs a pre-flight health check that verifies connectivity and
model availability, and then transfers control to the FSM. On
shutdown it ensures that all threads are joined cleanly and that the
Arduino is commanded to idle before the process exits. The Application
Layer deliberately contains no business logic — it only orchestrates
the lifecycle of the layers below it~\cite{bass2021software}.
 
\paragraph{5. Utility Layer.}
The Utility Layer provides shared infrastructure consumed by all other
layers: a structured logging system that tags every record with
timestamp, thread identity, and severity; a configuration manager
that loads parameters from a single file and exposes them through a
typed interface; and lightweight shared data structures (thread-safe
queues, event objects) used for inter-layer communication. Centralising
these concerns prevents duplication and ensures consistent behaviour
across the entire codebase~\cite{martin2008clean}.
 
% ............................................................
\subsection{Software Engineering Principles Applied}
% ............................................................
 
The architecture embodies several established software engineering
principles that improve long-term maintainability and testability:
 
\begin{itemize}
    \item \textbf{Separation of concerns.} Each layer and each service
          has a distinct responsibility. Perception, communication,
          control, and monitoring are never mixed within a single
          module~\cite{dijkstra1982role}.
 
    \item \textbf{Single responsibility principle.} Every class or
          module addresses exactly one aspect of system
          behaviour~\cite{martin2008clean}. This makes individual
          components easier to test in isolation and reduces the risk
          that a change in one area breaks another.
 
    \item \textbf{Dependency injection.} Services receive their
          dependencies (hardware drivers, configuration objects,
          shared queues) as constructor arguments rather than
          instantiating them internally. This decouples modules from
          concrete implementations and enables mock objects to be
          substituted during testing~\cite{gamma1994design}.
 
    \item \textbf{Fail-fast error handling.} Errors are detected and
          propagated upward to the FSM as early as possible rather
          than being silently ignored or retried indefinitely. This
          prevents the system from operating in an undefined
          degraded state~\cite{martin2008clean}.
 
    \item \textbf{Centralised configuration management.} All tunable
          parameters — timeouts, model paths, server addresses,
          camera resolution — are stored in a single configuration
          file managed by the Utility Layer. No magic numbers appear
          in application code.
\end{itemize}
 
Table~\ref{tab:sw_principles} maps each principle to the layer or
component where it is primarily enforced.
 
\begin{table}[h]
\centering
\caption{Software engineering principles and their primary enforcement points.}
\label{tab:sw_principles}
\begin{tabular}{ll}
\hline
\textbf{Principle} & \textbf{Primary Enforcement Point} \\
\hline
Separation of concerns      & All inter-layer boundaries          \\
Single responsibility       & Individual service classes          \\
Dependency injection        & Application Layer (object wiring)   \\
Fail-fast error handling    & FSM + Watchdog                      \\
Centralised configuration   & Utility Layer (config manager)      \\
\hline
\end{tabular}
\end{table}
 
\bigskip
\noindent\textbf{[DIAGRAM 4 — Startup and Shutdown Sequence Diagram]}\\
\textit{A UML sequence diagram showing the Application Layer
orchestrating startup: Config load $\rightarrow$ Hardware Layer init
$\rightarrow$ Services Layer init $\rightarrow$ Health check
$\rightarrow$ FSM transition to Operational. A second sequence
(right side or below) shows the shutdown path: Fault event
$\rightarrow$ FSM $\rightarrow$ Services teardown $\rightarrow$
Arduino idle command $\rightarrow$ process exit. Place after
Table~\ref{tab:sw_principles}.}
\bigskip

\section{Communication Architecture}
\label{sec:communication}
% ------------------------------------------------------------
 
The system employs a \emph{two-channel distributed communication
architecture} in which the Raspberry Pi acts as a central gateway,
bridging a network-facing channel on one side and a hardware-facing
channel on the other. Each channel is designed independently to match
the specific requirements of its domain — bandwidth, latency,
reliability, and message semantics — while the Raspberry Pi
translates between them, insulating network logic from hardware
control logic and vice versa~\cite{tanenbaum2021networks}.
 
\bigskip
\noindent\textbf{[DIAGRAM 1 — Communication Architecture Overview]}\\
\textit{A three-node horizontal diagram: Server (left)
$\xleftrightarrow{\text{WebSocket / TCP}}$ Raspberry Pi (centre)
$\xleftrightarrow{\text{UART / USB}}$ Arduino (right). Annotate each
arrow with message categories: outbound from Pi to server (state,
telemetry, detections), inbound from server to Pi (commands), outbound
from Pi to Arduino (actuator commands), inbound from Arduino to Pi
(ACKs, status). Place after this introduction paragraph.}
\bigskip
 
% ............................................................
\subsection{Server-to-Raspberry Pi Communication — WebSocket}
% ............................................................
 
The link between the backend server and the Raspberry Pi is
implemented using the \textbf{WebSocket protocol}
(RFC~6455)~\cite{fette2011websocket}. WebSocket upgrades a standard
HTTP connection into a persistent, full-duplex TCP channel: once the
handshake is complete, either party can push a message at any time
without waiting for a request from the other side. This property makes
WebSocket well-suited for real-time telemetry and control, where the
server must be able to issue commands reactively and the robot must be
able to report events as they occur, with minimal
latency~\cite{lubbers2011websocket}.
 
The Raspberry Pi establishes the connection as a client, initiating an
outbound connection to a dedicated control endpoint on the server.
Maintaining the Pi as the initiating party simplifies network
configuration — the server does not need to reach the Pi's private
address, and no inbound port-forwarding rules are required on the
robot's network.
 
\paragraph{Raspberry Pi $\to$ Server (uplink).}
The following message categories are sent by the Pi to the server:
 
\begin{itemize}
    \item \textbf{Robot state.} Current FSM state (Idle, Streaming,
          Detection, Actuation, Error), enabling the dashboard to
          reflect live operational status.
    \item \textbf{Periodic status updates.} Heartbeat frames confirming
          that the Pi process is alive and the connection is healthy.
    \item \textbf{System health telemetry.} CPU load, memory
          utilisation, core temperature, and uptime — collected by the
          Health Monitor service described in
          Section~\ref{sec:software}.
    \item \textbf{AI detection results.} Bounding box coordinates,
          class labels, and confidence scores produced by the
          Inference Service, forwarded to the server for logging and
          display on the dashboard.
\end{itemize}
 
\paragraph{Server $\to$ Raspberry Pi (downlink).}
The following command categories are received by the Pi from the
server:
 
\begin{itemize}
    \item \textbf{Movement commands.} Directional instructions
          (forward, reverse, turn left, turn right, stop) to be
          translated into motor commands and forwarded to the Arduino.
    \item \textbf{Operational mode changes.} Instructions to switch
          between local and remote inference modes, or to activate
          and deactivate the detection pipeline.
    \item \textbf{Inference control.} Start and stop commands for
          the AI processing loop, allowing the operator to enable
          detection only when needed to reduce computational load.
    \item \textbf{System management.} Shutdown and reboot commands
          that trigger a controlled teardown sequence through the FSM.
\end{itemize}
 
% ............................................................
\subsection{Raspberry Pi-to-Arduino Communication — UART Serial}
% ............................................................
 
The link between the Raspberry Pi and the Arduino is implemented as a
\textbf{UART serial connection over USB}. Unlike WebSocket, which
operates over a network stack with variable latency, UART provides a
point-to-point, hardware-level byte stream with deterministic
timing — an appropriate choice for delivering actuator commands that
must be executed promptly and reliably~\cite{axelson2007serial}.
 
The USB connection serves a dual purpose: it carries the serial data
stream \emph{and} provides 5\,V power to the Arduino, eliminating the
need for a separate power cable between the two boards (as noted in
Section~\ref{sec:hardware}).
 
\paragraph{Protocol design.}
The protocol follows a \emph{command-response model}: the Raspberry Pi
sends a command frame, and the Arduino replies with an acknowledgement
or status report before the Pi sends the next command. This simple
half-duplex exchange provides built-in flow control without a complex
protocol stack~\cite{axelson2007serial}.
 
Commands are encoded as \textbf{newline-terminated ASCII strings}
rather than a compact binary format. While binary encoding would reduce
per-message byte count, ASCII provides two practical advantages during
development and debugging: a serial monitor (e.g., the Arduino IDE
Serial Monitor) can display all traffic in human-readable form without
a protocol decoder, and malformed messages are immediately recognisable
by inspection.
 
\paragraph{Supported command categories.}
\begin{itemize}
    \item \textbf{Motor control.} Speed and direction commands for the
          four drive motors, mapped to PWM duty cycles by the
          Arduino's Hardware Abstraction Layer.
    \item \textbf{Pump control.} Activation and deactivation of the
          spraying pump mechanism.
    \item \textbf{Status request.} A polling command that causes the
          Arduino to report the current state of all actuators.
    \item \textbf{Emergency stop.} A priority command that immediately
          idles all actuators regardless of current state, mirroring
          the hardware-level failsafe described in
          Section~\ref{sec:software}.
\end{itemize}
 
\paragraph{Arduino responses.}
The Arduino replies to every command with one of three response types:
an \emph{acknowledgement} confirming successful execution, a
\emph{status report} enumerating actuator states, or an \emph{error
message} indicating a validation failure or hardware fault. This
feedback loop allows the Raspberry Pi's Communication Service to
detect silent failures — cases where a command was sent but no
response was received within the expected window — and escalate them
to the FSM as fault events.
 
% ............................................................
\subsection{End-to-End Data Flow}
% ............................................................
 
Figure~\ref{fig:dataflow} illustrates the complete data flow through
the communication architecture, from user input at the dashboard to
physical actuation at the robot, and from sensor data at the robot
back to the dashboard.
 
\bigskip
\noindent\textbf{[DIAGRAM 2 — End-to-End Data Flow]}\\
\textit{A vertical flow diagram with four horizontal tiers: Dashboard
(top), Server, Raspberry Pi, Arduino / Hardware (bottom). Downward
arrows carry control commands; upward arrows carry telemetry and
detection results. Label each arrow with the protocol and example
message type. The Raspberry Pi tier should show the internal
translation step (WebSocket message parsed $\to$ UART command
generated). Reference this as Figure~\ref{fig:dataflow} in the text.
Place immediately here.}
\label{fig:dataflow}
\bigskip
 
\paragraph{Control path (top-down).}
\begin{enumerate}
    \item The operator issues a command (e.g., ``move forward'') via
          the dashboard.
    \item The dashboard sends the command to the server over HTTP or
          WebSocket.
    \item The server forwards the command to the Raspberry Pi over the
          persistent WebSocket connection.
    \item The Raspberry Pi's Communication Service receives the message,
          validates it, and translates it into the corresponding UART
          command string.
    \item The command is transmitted to the Arduino over the serial
          link.
    \item The Arduino's Command Processor decodes the instruction and
          updates the Actuator Abstraction Layer, which drives the
          motor PWM signals accordingly.
\end{enumerate}
 
\paragraph{Data path (bottom-up).}
\begin{enumerate}
    \item The camera acquires a frame; the Inference Service runs
          detection and produces a result.
    \item The Arduino periodically reports actuator status to the
          Raspberry Pi via the serial acknowledgement channel.
    \item The Raspberry Pi aggregates detection results, actuator
          status, and health telemetry, then forwards them to the
          server over the WebSocket uplink.
    \item The server relays the data to the dashboard, which updates
          the live display in real time.
\end{enumerate}
 
\paragraph{Gateway role of the Raspberry Pi.}
The Pi's role as a \emph{communication gateway}~\cite{tanenbaum2021networks}
is central to the architecture's modularity. Because the server and
the Arduino never communicate directly, either channel can be modified
or replaced — for example, substituting MQTT for WebSocket on the
network side, or upgrading to CAN bus on the hardware side — without
requiring changes to the other. The only component that must be updated
is the Pi's Communication Service, which is isolated in its own
module precisely for this reason.
 
Table~\ref{tab:comms_compare} summarises and contrasts the two
communication channels.
 
\begin{table}[h]
\centering
\caption{Comparison of the two communication channels.}
\label{tab:comms_compare}
\begin{tabular}{lll}
\hline
\textbf{Property}       & \textbf{WebSocket (Server $\leftrightarrow$ Pi)}
                        & \textbf{UART (Pi $\leftrightarrow$ Arduino)} \\
\hline
Physical medium         & Wi-Fi / Ethernet (TCP/IP)  & USB cable           \\
Protocol standard       & RFC 6455                   & UART / USB-CDC      \\
Message format          & JSON / binary frames       & ASCII, newline-delimited \\
Directionality          & Full-duplex                & Half-duplex (command-response) \\
Latency characteristic  & Variable (network-bound)   & Deterministic       \\
Primary concern         & Bandwidth, remote reach    & Timing, reliability \\
Error detection         & TCP retransmission         & Application-level ACK \\
\hline
\end{tabular}
\end{table}
 
% ============================================================
%  Chapter 3 — Conclusion
% ============================================================
 
\section*{Chapter Summary}
\addcontentsline{toc}{section}{Chapter Summary}
 
This chapter presented the complete system design of the proposed
agricultural robotic platform across four interconnected dimensions.
 
The \textbf{high-level architecture} (Section~\ref{sec:architecture})
established the three-tier structure — embedded device, server, and
dashboard — and introduced the event-driven, multi-threaded execution
model that allows perception, communication, and control to proceed
concurrently without blocking one another.
 
The \textbf{hardware architecture} (Section~\ref{sec:hardware})
defined the physical component set: the Raspberry Pi as the central
processing unit, the Arduino as the real-time actuator controller,
the camera module for visual perception, and the motor and pump
subsystems for mobility and pesticide delivery. The two-domain power
supply strategy was described as essential for preventing
high-current actuator transients from corrupting the computational
supply rail. The boundary between the implemented prototype and the
full conceptual design was made explicit.
 
The \textbf{software architecture} (Section~\ref{sec:software})
described the layered design of both the Raspberry Pi stack — five
layers from hardware abstraction up to the application entry point —
and the Arduino firmware — four modules from command parsing down to
hardware signal generation. The Finite State Machine, watchdog
mechanisms, and fail-safe behaviours that make the system
self-supervising and fault-tolerant were examined in detail.
 
The \textbf{communication architecture} (Section~\ref{sec:communication})
explained the two-channel design: WebSocket for full-duplex,
low-latency network communication with the server, and UART serial
for deterministic, auditable command delivery to the Arduino. The
end-to-end data flow — from operator input at the dashboard through
to physical actuation, and from sensor data back to the live display
— was traced in both directions.
 
Together, these four sections establish that the system design is
modular, fault-tolerant, and extensible. The deliberate separation of
concerns at every layer — between network and hardware communication,
between high-level logic and real-time control, between perception and
actuation — ensures that each subsystem can be developed, tested, and
upgraded independently.
 
Chapter~4 proceeds to document the implementation of this design,
covering the specific tools, libraries, and engineering decisions made
during development, and the challenges encountered in translating the
architectural model into a working prototype.